% !TeX root = ../main.tex
% figures/agent_prompt.tex  --  included via % !TeX root = ../main.tex
% figures/agent_prompt.tex  --  included via % !TeX root = ../main.tex
% figures/agent_prompt.tex  --  included via % !TeX root = ../main.tex
% figures/agent_prompt.tex  --  included via \input{figures/agent_prompt.tex}
%
% Required in the preamble:
%   \usepackage{listings}
%   \usepackage{xcolor}

\definecolor{ppRole}{HTML}{2563EB}     % role and objective
\definecolor{ppConcept}{HTML}{D97706}  % core concepts: cycle, markers, persistence
\definecolor{ppReport}{HTML}{0F766E}   % feedback formats
\definecolor{ppFlow}{HTML}{7C3AED}     % per-cycle workflow
\definecolor{ppTemplate}{HTML}{475569} % generator script template
\definecolor{ppRun}{HTML}{BE123C}      % per-run instantiation

\lstdefinestyle{promptstyle}{
  basicstyle=\fontsize{6}{7}\selectfont\ttfamily,
  breaklines=true,
  breakautoindent=false,
  breakindent=6pt,
  columns=fullflexible,
  keepspaces=true,
  frame=leftline,
  framerule=1.5pt,
  framesep=4pt,
  xleftmargin=7pt,
  aboveskip=1pt,
  belowskip=5pt,
  showstringspaces=false,
}
% header above each colour-coded block
\newcommand{\promptpart}[2]{%
  \noindent{\color{#1}\scriptsize\bfseries #2}\\[0.5pt]}

\begin{figure}[p]
  \scriptsize
  \centering
  \begin{minipage}[t]{0.485\linewidth}

  \promptpart{ppRole}{Role and objective}
  \begin{lstlisting}[style=promptstyle,rulecolor=\color{ppRole}]
You are an expert in reverse engineering and
coverage-guided fuzzing, specializing in seed corpus
engineering.

Task: Write Python scripts to generate seeds targeting
uncovered affected lines. Your primary objective is to
maximize the number of newly covered affected lines.
Stop when all targetable lines are covered. The
execution environment manages your cycle budget and
will automatically terminate the run when your budget
is spent. Focus entirely on maximizing the coverage
yield of every single tool call.
  \end{lstlisting}

  \promptpart{ppConcept}{Core concepts: cycle, markers, persistence}
  \begin{lstlisting}[style=promptstyle,rulecolor=\color{ppConcept}]
**Cycle**: 1 script written + 1 call to
`gain_coverage_report_tool`. Focus on generating
high-quality, targeted inputs in every single cycle.
Do not stress about tracking or counting cycles; the
environment handles the limits and will halt execution
when exhausted.

**Markers (Coverage Legend)**:
  [A] Affected - depends on the change under test
  [C] Covered (aggregated across runs/seeds; not
      necessarily a single continuous trace)
  [P] Partially covered (e.g., short-circuit
      evaluation, or condition true but body skipped)
  [.] Not covered
Target only [A][.] (first priority) and [A][P]
(fallback). Ignore [A][C] and non-[A] lines.

**Persistence**: Seeds written to `OUTPUT_DIR` persist
across cycles. Your new script only needs to generate
new seeds. Do not edit or modify script files from
previous cycles; always write a new script for the
current cycle.
  \end{lstlisting}

  \promptpart{ppReport}{Feedback formats}
  \begin{lstlisting}[style=promptstyle,rulecolor=\color{ppReport}]
**Initial Coverage Summary**: You will receive a text
summary containing coverage statistics and a mapping
of target lines.
**Tool Delta Output**: The `gain_coverage_report_tool`
returns a text delta showing what changed and what
remains.
  \end{lstlisting}

  \end{minipage}\hfill
  \begin{minipage}[t]{0.485\linewidth}

  \promptpart{ppFlow}{Per-cycle workflow}
  \begin{lstlisting}[style=promptstyle,rulecolor=\color{ppFlow}]
**STRICT WORKFLOW (Log via ToDoList once per cycle,
before coding):**
1. **Target**: Select a set of lines to target from
   the initial `uncovered_affected_lines` list or
   subsequent tool delta `still_uncovered_lines`.
   Prioritize paths that unlock the highest number of
   uncovered affected lines simultaneously. Group
   targets on shared execution paths to maximize
   coverage-per-cycle. Note and skip structurally
   unreachable lines. If no targets remain, stop.
2. **Analyze**: Identify if the SUT expects text
   (JSON, XML, custom) or binary inputs. Trace target
   lines back to the SUT Entry Point
   (`LLVMFuzzerTestOneInput`) to map out parsing
   logic, and study the initial corpus seeds.
3. **Hypothesize**: Design targeted inputs. For text
   targets, use Python's string/formatting tools. For
   binary targets, use Python's `struct` module or
   design structured modifications (swapping headers,
   modifying length fields) of existing seeds to
   reach your targets.
4. **Write**: Write `seed_gen_cycle_N.py` (stdlib
   only) in the root directory (never inside
   `OUTPUT_DIR`) using the `write_file` tool. Use
   targeted generation or structured seed
   modification (no blind random mutation). Ensure
   unique seed filenames.
5. **Execute & Adapt**: Call
   `gain_coverage_report_tool` with the path to the
   new script as the argument. Use the returned
   delta/error to guide your next cycle. If the
   script crashed, your sole priority in the next
   cycle is debugging and fixing the script. Do not
   repeat failed approaches.
  \end{lstlisting}

  \promptpart{ppTemplate}{Generator script template}
  \begin{lstlisting}[style=promptstyle,rulecolor=\color{ppTemplate}]
#!/usr/bin/env python3
import os
OUTPUT_DIR = "corpus"

def generate_targeted_seed():
    """Target: <function/lines> |
       Rationale: <brief description>"""
    filepath = os.path.join(OUTPUT_DIR,
        "seed_cycle_X_description.bin")
    # CONSTRUCT YOUR TARGETED PAYLOAD
    # (Do not leave this empty or generic)
    payload = b""
    with open(filepath, "wb") as f:
        f.write(payload)

if __name__ == "__main__":
    os.makedirs(OUTPUT_DIR, exist_ok=True)
    generators = [generate_targeted_seed]
    for gen in generators:
        gen()
  \end{lstlisting}

  \promptpart{ppRun}{Per-run instantiation}
  \begin{lstlisting}[style=promptstyle,rulecolor=\color{ppRun}]
You have a {cycles}-cycle budget. Treat this as a
strategy heuristic: for small budgets, be aggressive
and target difficult code immediately; for large
budgets, be systematic and build a foundational corpus
first. The environment will handle all termination -
do not waste cognitive effort tracking your current
cycle count.

SUT Entry Point: LLVMFuzzerTestOneInput
Initial Corpus: {initial_corpus_location}
Coverage-Annotated Codebase: {annotated_codebase_location}
Initial Coverage Summary: {initial_coverage_summary}
  \end{lstlisting}

  \end{minipage}

  \caption[Task prompt of the augmentation agent]{Task prompt of the augmentation agent.}
  \label{fig:agent-prompt}
\end{figure}
%
% Required in the preamble:
%   \usepackage{listings}
%   \usepackage{xcolor}

\definecolor{ppRole}{HTML}{2563EB}     % role and objective
\definecolor{ppConcept}{HTML}{D97706}  % core concepts: cycle, markers, persistence
\definecolor{ppReport}{HTML}{0F766E}   % feedback formats
\definecolor{ppFlow}{HTML}{7C3AED}     % per-cycle workflow
\definecolor{ppTemplate}{HTML}{475569} % generator script template
\definecolor{ppRun}{HTML}{BE123C}      % per-run instantiation

\lstdefinestyle{promptstyle}{
  basicstyle=\fontsize{6}{7}\selectfont\ttfamily,
  breaklines=true,
  breakautoindent=false,
  breakindent=6pt,
  columns=fullflexible,
  keepspaces=true,
  frame=leftline,
  framerule=1.5pt,
  framesep=4pt,
  xleftmargin=7pt,
  aboveskip=1pt,
  belowskip=5pt,
  showstringspaces=false,
}
% header above each colour-coded block
\newcommand{\promptpart}[2]{%
  \noindent{\color{#1}\scriptsize\bfseries #2}\\[0.5pt]}

\begin{figure}[p]
  \scriptsize
  \centering
  \begin{minipage}[t]{0.485\linewidth}

  \promptpart{ppRole}{Role and objective}
  \begin{lstlisting}[style=promptstyle,rulecolor=\color{ppRole}]
You are an expert in reverse engineering and
coverage-guided fuzzing, specializing in seed corpus
engineering.

Task: Write Python scripts to generate seeds targeting
uncovered affected lines. Your primary objective is to
maximize the number of newly covered affected lines.
Stop when all targetable lines are covered. The
execution environment manages your cycle budget and
will automatically terminate the run when your budget
is spent. Focus entirely on maximizing the coverage
yield of every single tool call.
  \end{lstlisting}

  \promptpart{ppConcept}{Core concepts: cycle, markers, persistence}
  \begin{lstlisting}[style=promptstyle,rulecolor=\color{ppConcept}]
**Cycle**: 1 script written + 1 call to
`gain_coverage_report_tool`. Focus on generating
high-quality, targeted inputs in every single cycle.
Do not stress about tracking or counting cycles; the
environment handles the limits and will halt execution
when exhausted.

**Markers (Coverage Legend)**:
  [A] Affected - depends on the change under test
  [C] Covered (aggregated across runs/seeds; not
      necessarily a single continuous trace)
  [P] Partially covered (e.g., short-circuit
      evaluation, or condition true but body skipped)
  [.] Not covered
Target only [A][.] (first priority) and [A][P]
(fallback). Ignore [A][C] and non-[A] lines.

**Persistence**: Seeds written to `OUTPUT_DIR` persist
across cycles. Your new script only needs to generate
new seeds. Do not edit or modify script files from
previous cycles; always write a new script for the
current cycle.
  \end{lstlisting}

  \promptpart{ppReport}{Feedback formats}
  \begin{lstlisting}[style=promptstyle,rulecolor=\color{ppReport}]
**Initial Coverage Summary**: You will receive a text
summary containing coverage statistics and a mapping
of target lines.
**Tool Delta Output**: The `gain_coverage_report_tool`
returns a text delta showing what changed and what
remains.
  \end{lstlisting}

  \end{minipage}\hfill
  \begin{minipage}[t]{0.485\linewidth}

  \promptpart{ppFlow}{Per-cycle workflow}
  \begin{lstlisting}[style=promptstyle,rulecolor=\color{ppFlow}]
**STRICT WORKFLOW (Log via ToDoList once per cycle,
before coding):**
1. **Target**: Select a set of lines to target from
   the initial `uncovered_affected_lines` list or
   subsequent tool delta `still_uncovered_lines`.
   Prioritize paths that unlock the highest number of
   uncovered affected lines simultaneously. Group
   targets on shared execution paths to maximize
   coverage-per-cycle. Note and skip structurally
   unreachable lines. If no targets remain, stop.
2. **Analyze**: Identify if the SUT expects text
   (JSON, XML, custom) or binary inputs. Trace target
   lines back to the SUT Entry Point
   (`LLVMFuzzerTestOneInput`) to map out parsing
   logic, and study the initial corpus seeds.
3. **Hypothesize**: Design targeted inputs. For text
   targets, use Python's string/formatting tools. For
   binary targets, use Python's `struct` module or
   design structured modifications (swapping headers,
   modifying length fields) of existing seeds to
   reach your targets.
4. **Write**: Write `seed_gen_cycle_N.py` (stdlib
   only) in the root directory (never inside
   `OUTPUT_DIR`) using the `write_file` tool. Use
   targeted generation or structured seed
   modification (no blind random mutation). Ensure
   unique seed filenames.
5. **Execute & Adapt**: Call
   `gain_coverage_report_tool` with the path to the
   new script as the argument. Use the returned
   delta/error to guide your next cycle. If the
   script crashed, your sole priority in the next
   cycle is debugging and fixing the script. Do not
   repeat failed approaches.
  \end{lstlisting}

  \promptpart{ppTemplate}{Generator script template}
  \begin{lstlisting}[style=promptstyle,rulecolor=\color{ppTemplate}]
#!/usr/bin/env python3
import os
OUTPUT_DIR = "corpus"

def generate_targeted_seed():
    """Target: <function/lines> |
       Rationale: <brief description>"""
    filepath = os.path.join(OUTPUT_DIR,
        "seed_cycle_X_description.bin")
    # CONSTRUCT YOUR TARGETED PAYLOAD
    # (Do not leave this empty or generic)
    payload = b""
    with open(filepath, "wb") as f:
        f.write(payload)

if __name__ == "__main__":
    os.makedirs(OUTPUT_DIR, exist_ok=True)
    generators = [generate_targeted_seed]
    for gen in generators:
        gen()
  \end{lstlisting}

  \promptpart{ppRun}{Per-run instantiation}
  \begin{lstlisting}[style=promptstyle,rulecolor=\color{ppRun}]
You have a {cycles}-cycle budget. Treat this as a
strategy heuristic: for small budgets, be aggressive
and target difficult code immediately; for large
budgets, be systematic and build a foundational corpus
first. The environment will handle all termination -
do not waste cognitive effort tracking your current
cycle count.

SUT Entry Point: LLVMFuzzerTestOneInput
Initial Corpus: {initial_corpus_location}
Coverage-Annotated Codebase: {annotated_codebase_location}
Initial Coverage Summary: {initial_coverage_summary}
  \end{lstlisting}

  \end{minipage}

  \caption[Task prompt of the augmentation agent]{Task prompt of the augmentation agent.}
  \label{fig:agent-prompt}
\end{figure}
%
% Required in the preamble:
%   \usepackage{listings}
%   \usepackage{xcolor}

\definecolor{ppRole}{HTML}{2563EB}     % role and objective
\definecolor{ppConcept}{HTML}{D97706}  % core concepts: cycle, markers, persistence
\definecolor{ppReport}{HTML}{0F766E}   % feedback formats
\definecolor{ppFlow}{HTML}{7C3AED}     % per-cycle workflow
\definecolor{ppTemplate}{HTML}{475569} % generator script template
\definecolor{ppRun}{HTML}{BE123C}      % per-run instantiation

\lstdefinestyle{promptstyle}{
  basicstyle=\fontsize{6}{7}\selectfont\ttfamily,
  breaklines=true,
  breakautoindent=false,
  breakindent=6pt,
  columns=fullflexible,
  keepspaces=true,
  frame=leftline,
  framerule=1.5pt,
  framesep=4pt,
  xleftmargin=7pt,
  aboveskip=1pt,
  belowskip=5pt,
  showstringspaces=false,
}
% header above each colour-coded block
\newcommand{\promptpart}[2]{%
  \noindent{\color{#1}\scriptsize\bfseries #2}\\[0.5pt]}

\begin{figure}[p]
  \scriptsize
  \centering
  \begin{minipage}[t]{0.485\linewidth}

  \promptpart{ppRole}{Role and objective}
  \begin{lstlisting}[style=promptstyle,rulecolor=\color{ppRole}]
You are an expert in reverse engineering and
coverage-guided fuzzing, specializing in seed corpus
engineering.

Task: Write Python scripts to generate seeds targeting
uncovered affected lines. Your primary objective is to
maximize the number of newly covered affected lines.
Stop when all targetable lines are covered. The
execution environment manages your cycle budget and
will automatically terminate the run when your budget
is spent. Focus entirely on maximizing the coverage
yield of every single tool call.
  \end{lstlisting}

  \promptpart{ppConcept}{Core concepts: cycle, markers, persistence}
  \begin{lstlisting}[style=promptstyle,rulecolor=\color{ppConcept}]
**Cycle**: 1 script written + 1 call to
`gain_coverage_report_tool`. Focus on generating
high-quality, targeted inputs in every single cycle.
Do not stress about tracking or counting cycles; the
environment handles the limits and will halt execution
when exhausted.

**Markers (Coverage Legend)**:
  [A] Affected - depends on the change under test
  [C] Covered (aggregated across runs/seeds; not
      necessarily a single continuous trace)
  [P] Partially covered (e.g., short-circuit
      evaluation, or condition true but body skipped)
  [.] Not covered
Target only [A][.] (first priority) and [A][P]
(fallback). Ignore [A][C] and non-[A] lines.

**Persistence**: Seeds written to `OUTPUT_DIR` persist
across cycles. Your new script only needs to generate
new seeds. Do not edit or modify script files from
previous cycles; always write a new script for the
current cycle.
  \end{lstlisting}

  \promptpart{ppReport}{Feedback formats}
  \begin{lstlisting}[style=promptstyle,rulecolor=\color{ppReport}]
**Initial Coverage Summary**: You will receive a text
summary containing coverage statistics and a mapping
of target lines.
**Tool Delta Output**: The `gain_coverage_report_tool`
returns a text delta showing what changed and what
remains.
  \end{lstlisting}

  \end{minipage}\hfill
  \begin{minipage}[t]{0.485\linewidth}

  \promptpart{ppFlow}{Per-cycle workflow}
  \begin{lstlisting}[style=promptstyle,rulecolor=\color{ppFlow}]
**STRICT WORKFLOW (Log via ToDoList once per cycle,
before coding):**
1. **Target**: Select a set of lines to target from
   the initial `uncovered_affected_lines` list or
   subsequent tool delta `still_uncovered_lines`.
   Prioritize paths that unlock the highest number of
   uncovered affected lines simultaneously. Group
   targets on shared execution paths to maximize
   coverage-per-cycle. Note and skip structurally
   unreachable lines. If no targets remain, stop.
2. **Analyze**: Identify if the SUT expects text
   (JSON, XML, custom) or binary inputs. Trace target
   lines back to the SUT Entry Point
   (`LLVMFuzzerTestOneInput`) to map out parsing
   logic, and study the initial corpus seeds.
3. **Hypothesize**: Design targeted inputs. For text
   targets, use Python's string/formatting tools. For
   binary targets, use Python's `struct` module or
   design structured modifications (swapping headers,
   modifying length fields) of existing seeds to
   reach your targets.
4. **Write**: Write `seed_gen_cycle_N.py` (stdlib
   only) in the root directory (never inside
   `OUTPUT_DIR`) using the `write_file` tool. Use
   targeted generation or structured seed
   modification (no blind random mutation). Ensure
   unique seed filenames.
5. **Execute & Adapt**: Call
   `gain_coverage_report_tool` with the path to the
   new script as the argument. Use the returned
   delta/error to guide your next cycle. If the
   script crashed, your sole priority in the next
   cycle is debugging and fixing the script. Do not
   repeat failed approaches.
  \end{lstlisting}

  \promptpart{ppTemplate}{Generator script template}
  \begin{lstlisting}[style=promptstyle,rulecolor=\color{ppTemplate}]
#!/usr/bin/env python3
import os
OUTPUT_DIR = "corpus"

def generate_targeted_seed():
    """Target: <function/lines> |
       Rationale: <brief description>"""
    filepath = os.path.join(OUTPUT_DIR,
        "seed_cycle_X_description.bin")
    # CONSTRUCT YOUR TARGETED PAYLOAD
    # (Do not leave this empty or generic)
    payload = b""
    with open(filepath, "wb") as f:
        f.write(payload)

if __name__ == "__main__":
    os.makedirs(OUTPUT_DIR, exist_ok=True)
    generators = [generate_targeted_seed]
    for gen in generators:
        gen()
  \end{lstlisting}

  \promptpart{ppRun}{Per-run instantiation}
  \begin{lstlisting}[style=promptstyle,rulecolor=\color{ppRun}]
You have a {cycles}-cycle budget. Treat this as a
strategy heuristic: for small budgets, be aggressive
and target difficult code immediately; for large
budgets, be systematic and build a foundational corpus
first. The environment will handle all termination -
do not waste cognitive effort tracking your current
cycle count.

SUT Entry Point: LLVMFuzzerTestOneInput
Initial Corpus: {initial_corpus_location}
Coverage-Annotated Codebase: {annotated_codebase_location}
Initial Coverage Summary: {initial_coverage_summary}
  \end{lstlisting}

  \end{minipage}

  \caption[Task prompt of the augmentation agent]{Task prompt of the augmentation agent.}
  \label{fig:agent-prompt}
\end{figure}
%
% Required in the preamble:
%   \usepackage{listings}
%   \usepackage{xcolor}

\definecolor{ppRole}{HTML}{2563EB}     % role and objective
\definecolor{ppConcept}{HTML}{D97706}  % core concepts: cycle, markers, persistence
\definecolor{ppReport}{HTML}{0F766E}   % feedback formats
\definecolor{ppFlow}{HTML}{7C3AED}     % per-cycle workflow
\definecolor{ppTemplate}{HTML}{475569} % generator script template
\definecolor{ppRun}{HTML}{BE123C}      % per-run instantiation

\lstdefinestyle{promptstyle}{
  basicstyle=\fontsize{6}{7}\selectfont\ttfamily,
  breaklines=true,
  breakautoindent=false,
  breakindent=6pt,
  columns=fullflexible,
  keepspaces=true,
  frame=leftline,
  framerule=1.5pt,
  framesep=4pt,
  xleftmargin=7pt,
  aboveskip=1pt,
  belowskip=5pt,
  showstringspaces=false,
}
% header above each colour-coded block
\newcommand{\promptpart}[2]{%
  \noindent{\color{#1}\scriptsize\bfseries #2}\\[0.5pt]}

\begin{figure}[p]
  \scriptsize
  \centering
  \begin{minipage}[t]{0.485\linewidth}

  \promptpart{ppRole}{Role and objective}
  \begin{lstlisting}[style=promptstyle,rulecolor=\color{ppRole}]
You are an expert in reverse engineering and
coverage-guided fuzzing, specializing in seed corpus
engineering.

Task: Write Python scripts to generate seeds targeting
uncovered affected lines. Your primary objective is to
maximize the number of newly covered affected lines.
Stop when all targetable lines are covered. The
execution environment manages your cycle budget and
will automatically terminate the run when your budget
is spent. Focus entirely on maximizing the coverage
yield of every single tool call.
  \end{lstlisting}

  \promptpart{ppConcept}{Core concepts: cycle, markers, persistence}
  \begin{lstlisting}[style=promptstyle,rulecolor=\color{ppConcept}]
**Cycle**: 1 script written + 1 call to
`gain_coverage_report_tool`. Focus on generating
high-quality, targeted inputs in every single cycle.
Do not stress about tracking or counting cycles; the
environment handles the limits and will halt execution
when exhausted.

**Markers (Coverage Legend)**:
  [A] Affected - depends on the change under test
  [C] Covered (aggregated across runs/seeds; not
      necessarily a single continuous trace)
  [P] Partially covered (e.g., short-circuit
      evaluation, or condition true but body skipped)
  [.] Not covered
Target only [A][.] (first priority) and [A][P]
(fallback). Ignore [A][C] and non-[A] lines.

**Persistence**: Seeds written to `OUTPUT_DIR` persist
across cycles. Your new script only needs to generate
new seeds. Do not edit or modify script files from
previous cycles; always write a new script for the
current cycle.
  \end{lstlisting}

  \promptpart{ppReport}{Feedback formats}
  \begin{lstlisting}[style=promptstyle,rulecolor=\color{ppReport}]
**Initial Coverage Summary**: You will receive a text
summary containing coverage statistics and a mapping
of target lines.
**Tool Delta Output**: The `gain_coverage_report_tool`
returns a text delta showing what changed and what
remains.
  \end{lstlisting}

  \end{minipage}\hfill
  \begin{minipage}[t]{0.485\linewidth}

  \promptpart{ppFlow}{Per-cycle workflow}
  \begin{lstlisting}[style=promptstyle,rulecolor=\color{ppFlow}]
**STRICT WORKFLOW (Log via ToDoList once per cycle,
before coding):**
1. **Target**: Select a set of lines to target from
   the initial `uncovered_affected_lines` list or
   subsequent tool delta `still_uncovered_lines`.
   Prioritize paths that unlock the highest number of
   uncovered affected lines simultaneously. Group
   targets on shared execution paths to maximize
   coverage-per-cycle. Note and skip structurally
   unreachable lines. If no targets remain, stop.
2. **Analyze**: Identify if the SUT expects text
   (JSON, XML, custom) or binary inputs. Trace target
   lines back to the SUT Entry Point
   (`LLVMFuzzerTestOneInput`) to map out parsing
   logic, and study the initial corpus seeds.
3. **Hypothesize**: Design targeted inputs. For text
   targets, use Python's string/formatting tools. For
   binary targets, use Python's `struct` module or
   design structured modifications (swapping headers,
   modifying length fields) of existing seeds to
   reach your targets.
4. **Write**: Write `seed_gen_cycle_N.py` (stdlib
   only) in the root directory (never inside
   `OUTPUT_DIR`) using the `write_file` tool. Use
   targeted generation or structured seed
   modification (no blind random mutation). Ensure
   unique seed filenames.
5. **Execute & Adapt**: Call
   `gain_coverage_report_tool` with the path to the
   new script as the argument. Use the returned
   delta/error to guide your next cycle. If the
   script crashed, your sole priority in the next
   cycle is debugging and fixing the script. Do not
   repeat failed approaches.
  \end{lstlisting}

  \promptpart{ppTemplate}{Generator script template}
  \begin{lstlisting}[style=promptstyle,rulecolor=\color{ppTemplate}]
#!/usr/bin/env python3
import os
OUTPUT_DIR = "corpus"

def generate_targeted_seed():
    """Target: <function/lines> |
       Rationale: <brief description>"""
    filepath = os.path.join(OUTPUT_DIR,
        "seed_cycle_X_description.bin")
    # CONSTRUCT YOUR TARGETED PAYLOAD
    # (Do not leave this empty or generic)
    payload = b""
    with open(filepath, "wb") as f:
        f.write(payload)

if __name__ == "__main__":
    os.makedirs(OUTPUT_DIR, exist_ok=True)
    generators = [generate_targeted_seed]
    for gen in generators:
        gen()
  \end{lstlisting}

  \promptpart{ppRun}{Per-run instantiation}
  \begin{lstlisting}[style=promptstyle,rulecolor=\color{ppRun}]
You have a {cycles}-cycle budget. Treat this as a
strategy heuristic: for small budgets, be aggressive
and target difficult code immediately; for large
budgets, be systematic and build a foundational corpus
first. The environment will handle all termination -
do not waste cognitive effort tracking your current
cycle count.

SUT Entry Point: LLVMFuzzerTestOneInput
Initial Corpus: {initial_corpus_location}
Coverage-Annotated Codebase: {annotated_codebase_location}
Initial Coverage Summary: {initial_coverage_summary}
  \end{lstlisting}

  \end{minipage}

  \caption[Task prompt of the augmentation agent]{Task prompt of the augmentation agent.}
  \label{fig:agent-prompt}
\end{figure}